\subsection{What the Solver Actually Forms}
\label{subsec:kkt}

Section~\ref{subsec:ocp} states the model in the solver's generic objects
$x, u, \ell, f, c$.  This subsection eliminates them: every quantity below is
written in the variables of \eqref{eq:cp_all} and the multipliers attached to
its constraints.  All expressions are verified against the derivative functions
the solver actually builds.

\subsubsection{Variables carried at stage $t$}
The state is the stored energy, and nothing else:
\begin{equation}
\underbrace{B^t}_{\text{state}}; \qquad
\underbrace{P_{\text{Subs}}^t,\; P_B^t,\; s^t}_{\text{controls}} .
\label{eq:kkt_vars}
\end{equation}
Each stage carries one multiplier per constraint:
\begin{alignat}{3}
&\varphi_{\text{bal}}^t   &\;\;:\;\;& P_{\text{Subs}}^t + P_B^t - p_L^t = 0
\label{eq:kkt_mult_bal}\\
&\varphi_{\text{soc}}^t    && \bigl(B^t - P_B^t \Delta t\bigr) - B^{\min} - s^t = 0
\label{eq:kkt_mult_soc}\\
&\nu_{\text{Subs}}^t       && P_{\text{Subs}}^t \geq 0
\label{eq:kkt_mult_ps}\\
&\nu_{\min}^t,\; \nu_{\max}^t && P_B^{\min} \leq P_B^t \leq P_B^{\max}
\label{eq:kkt_mult_pb}\\
&\lambda_{B\min}^t,\; \lambda_{B\max}^t && 0 \leq s^t \leq B^{\max} - B^{\min} .
\label{eq:kkt_mult_b}
\end{alignat}
The last line is the correspondence that matters: because $s^t$ is defined by
\eqref{eq:kkt_mult_soc} as the slack of $B^{t+1}$ above $B^{\min}$, its two box
multipliers \emph{are} the energy-bound multipliers,
$\lambda_{B\min}^t$ on $B^{t+1} \geq B^{\min}$ and $\lambda_{B\max}^t$ on
$B^{t+1} \leq B^{\max}$.  There is no separate state-constraint dual.
$P_{\text{Subs}}^t$ has no upper bound, so no multiplier is attached to one.

\subsubsection{Stage derivatives}
With the dynamics $B^{t+1} = B^t - P_B^t \Delta t$ linear and the constraints
\eqref{eq:kkt_mult_bal}--\eqref{eq:kkt_mult_soc} linear in the controls, every
second derivative of $f$ and of $c^t$ vanishes, as does the cross term
$\partial^2 \ell^t / \partial u\, \partial B$.  What
remains is
\begin{alignat}{2}
&\frac{\partial c}{\partial(P_{\text{Subs}}, P_B, s)}
  &&= \begin{bmatrix} 1 & 1 & 0 \\ 0 & -\Delta t & -1 \end{bmatrix}
\label{eq:kkt_cu} \\
&\frac{\partial c}{\partial B}
  &&= \begin{bmatrix} 0 \\ 1 \end{bmatrix}
\label{eq:kkt_cx} \\
&\frac{\partial B^{t+1}}{\partial P_B^t} &&= -\Delta t
\label{eq:kkt_fu} \\
&\frac{\partial^2 \ell}{\partial (P_{\text{Subs}}, P_B, s)^2}
  &&= \operatorname{diag}\bigl(0,\; 2 C_B \Delta t,\; 0\bigr) .
\label{eq:kkt_luu}
\end{alignat}
The terminal stage adds $2w$ to the $B$ curvature, $2w(\Delta t)^2$ to the
$P_B$ curvature, and $-2w\Delta t$ to their cross term.

\subsubsection{The matrix that is factorised}
Write $V_{BB}^{t+1}$ for the curvature of the cost-to-go in the stored energy.
Collecting \eqref{eq:kkt_fu}--\eqref{eq:kkt_luu} with the barrier terms, the
stage Hessian is \emph{diagonal}:
\begin{equation}
\hat{H}^t = \operatorname{diag}\Bigl(
  \underbrace{\tfrac{\nu_{\text{Subs}}^t}{P_{\text{Subs}}^t}}_{P_{\text{Subs}}},\;\;
  \underbrace{2 C_B \Delta t + (\Delta t)^2 V_{BB}^{t+1} + \Sigma_{P_B}^t}_{P_B},\;\;
  \underbrace{\Sigma_{s}^t}_{s}
\Bigr),
\label{eq:kkt_H}
\end{equation}
\begin{align}
\Sigma_{P_B}^t &= \frac{\nu_{\min}^t}{P_B^t - P_B^{\min}}
                + \frac{\nu_{\max}^t}{P_B^{\max} - P_B^t},
\label{eq:kkt_sigma_pb} \\
\Sigma_{s}^t   &= \frac{\lambda_{B\min}^t}{s^t}
                + \frac{\lambda_{B\max}^t}{B^{\max} - B^{\min} - s^t} .
\label{eq:kkt_sigma_s}
\end{align}
Three controls and two equalities leave a \emph{one}-dimensional null space, so
after eliminating \eqref{eq:kkt_cu} exactly one direction survives:
\begin{equation}
d = \bigl(\Delta P_{\text{Subs}},\, \Delta P_B,\, \Delta s\bigr)
  = \bigl(-1,\; +1,\; -\Delta t\bigr) ,
\label{eq:kkt_null}
\end{equation}
which is the statement that the only freedom left at a stage is to move power
between the substation and the battery, the energy slack following.  The
quantity the solver factorises is therefore the \emph{scalar}
\begin{equation}
d^\top \hat{H}^t d
 = \frac{\nu_{\text{Subs}}^t}{P_{\text{Subs}}^t}
 + 2 C_B \Delta t + (\Delta t)^2 V_{BB}^{t+1}
 + \Sigma_{P_B}^t + (\Delta t)^2\, \Sigma_{s}^t ,
\label{eq:kkt_reduced}
\end{equation}
and the Cholesky whose failure triggers regularisation is a positivity test on
this one number.  Dropping the energy bound removes $s^t$ and its two
multipliers, leaving $d = (-1, +1)$ and
$d^\top \hat{H}^t d = \nu_{\text{Subs}}^t / P_{\text{Subs}}^t + 2 C_B \Delta t
+ (\Delta t)^2 V_{BB}^{t+1} + \Sigma_{P_B}^t$; the base instance strips the
barrier terms too, leaving $2 C_B \Delta t + (\Delta t)^2 V_{BB}^{t+1}$, which is
where $C_B$ enters as the curvature that keeps the problem solvable.

\subsubsection{What is passed backward}
Let $V_B^t$ denote the gradient of the cost-to-go in the stored energy.  With a
one-dimensional state, $V_B^t$ and $V_{BB}^t$ are scalars, so the whole backward
recursion is scalar arithmetic.  The
backward recursion carried between stages is
\begin{equation}
V_B^{t} = V_B^{t+1} + \varphi_{\text{soc}}^t
        + \bigl(\beta^\top \hat{Q}_u^t\bigr)_B + \bigl(\omega^\top c^t\bigr)_B ,
\label{eq:kkt_Vx}
\end{equation}
\begin{equation}
V_{BB}^{t} = V_{BB}^{t+1} + \bigl(\beta^\top \hat{B}^t\bigr)_{BB}
           + \bigl(\omega^\top c_x^t\bigr)_{BB} ,
\label{eq:kkt_Vxx}
\end{equation}
where $\beta$ and $\omega$ are the feedback gains produced by the stage solve.
At a feasible iterate ($c^t = 0$) that has converged ($\hat{Q}_u^t = 0$),
\eqref{eq:kkt_Vx} collapses to
\begin{equation}
V_B^{t} - V_B^{t+1} = \varphi_{\text{soc}}^t ,
\label{eq:kkt_costate}
\end{equation}
i.e.\ the value gradient in stored energy is the discrete costate, and its
stage-to-stage increment is exactly the multiplier on the energy relation.
Equation~\eqref{eq:kkt_Vxx} has no counterpart in a first-order method: it is the
curvature, and it is the object that a scheme passing only
\eqref{eq:kkt_costate} between stages does not carry.

\subsubsection{What is passed forward}
With step size $\gamma$ and $\delta B = B^t - \bar{B}^t$ the deviation from the
nominal trajectory, every quantity is updated by the same affine rule:
\begin{alignat}{2}
P_{\text{Subs}}^t &= \bar{P}_{\text{Subs}}^t + \gamma\,\alpha_{\text{Subs}}^t
   + \beta_{\text{Subs}}^t\, \delta B , \quad&
P_B^t &= \bar{P}_B^t + \gamma\,\alpha_{B}^t + \beta_{B}^t\, \delta B ,
\label{eq:kkt_fwd_u} \\
\varphi_\bullet^t &= \bar{\varphi}_\bullet^t + \gamma\,\psi_\bullet^t
   + \omega_\bullet^t\, \delta B , &
\nu_\bullet^t &= \bar{\nu}_\bullet^t + \gamma\,\chi_\bullet^t
   + \zeta_\bullet^t\, \delta B ,
\label{eq:kkt_fwd_mult}
\end{alignat}
so the primal variables, the equality multipliers and the bound multipliers are
all propagated together along the same search direction, with a single $\gamma$
accepted by the filter for all of them.
